% DEFINED:   sec:time-equivariance, sec:bptt, sec:vanishing-gradients,
%            eq:rnn-jacobian
% RESOLVED:  ch:cnn, sec:permuted-pixel, sec:mlp-backprop (backward into
%            Part II and Part I)
% FORWARD:   (none; the experiment and the Transformer are prose-only)

\section{Time-translation equivariance}\label{sec:time-equivariance}

The argument from \Cref{ch:cnn} transfers without change. Replace shift in
space with shift in time, and every line still holds.

Suppose we shift the input sequence by $\Delta t$, so the new sequence is
$x'_t = x_{t + \Delta t}$. Starting from the same initial state, the new
sequence of hidden states is $h'_t = h_{t + \Delta t}$. The cell applies one
computation at every step, so shifting the input shifts the output by the
same amount. This is a property of the operator, not of any particular
weights: it holds at initialization, it holds after training, and it holds
for every weight setting in between. Weight sharing across time is
time-translation equivariance, in the same sense that weight sharing across
space was translation equivariance for images.

(Strictly, the equivariance holds when the initial state $h_0$ is the same.
In practice we reset $h_0 = 0$ between sequences, which breaks the symmetry
at sequence boundaries. The within-sequence computation is fully
time-translation-equivariant, and that is what the inductive-bias claim
rests on.)

Equivariance is exactly the property the permuted-pixel control of
\Cref{sec:permuted-pixel} probed, now along the time axis. There the test
destroyed spatial structure by applying one fixed permutation to the pixels
of every image. The convolutional network, which assumes neighbors are
neighbors, lost more accuracy than the multilayer perceptron, which assumes
nothing about layout, because the convolutional network was actually using
the locality its weights bake in. The analogue here is a multilayer
perceptron on a flattened sequence. It would learn the phrase ``the cat'' at
positions one through three and then fail to recognize the same phrase at
positions five through seven, because it holds separate weights for each
position and the two occurrences share none. The recurrent cell recognizes
the phrase wherever it falls, because the same weights read every position.
That is the prior earning its keep, and it is the same prior the
permuted-pixel control of \Cref{ch:cnn} confirmed for space.

\section{Unrolling and backprop through time}\label{sec:bptt}

Training a recurrent network means computing the gradient of the loss with
respect to $W_{xh}$, $W_{hh}$, and $b_h$ for an unrolled sequence of length
$T$. The standard technique is backpropagation through time. Unroll the cell
$T$ times, treat the result as a deep feedforward network with $T$ layers
that happen to share one set of weights, and apply the usual chain rule.

The forward pass processes $x_1, \ldots, x_T$ in order, threading the state
$h_0 \to h_1 \to \cdots \to h_T$. At each step an output layer, here a
\texttt{Linear} projection to vocabulary logits followed by
\texttt{SoftmaxCrossEntropy}, contributes a per-step loss $L_t$. The total
loss is the sum over steps:
\[
  L = \sum_{t=1}^{T} L_t.
\]

The backward pass is where the recurrence shows up. The gradient on $h_t$
arrives from two places:
\[
  \frac{\partial L}{\partial h_t}
  =
  \underbrace{\frac{\partial L_t}{\partial h_t}}_{\text{from $L_t$ directly}}
  +
  \underbrace{\frac{\partial L}{\partial h_{t+1}}\,
              \frac{\partial h_{t+1}}{\partial h_t}}_{\text{from the future,
              through $h_{t+1}$}}.
\]
The first term is the gradient that the output projection at step $t$ sends
back. The second is everything downstream of step $t$, funneled through the
single channel the state provides, $h_t$ feeding $h_{t+1}$. In code this is
one extra argument on the cell's \texttt{backward}. On top of
\texttt{dh\_out} from the output projection, the cell receives
\texttt{dstate\_next}, the gradient on $h_t$ that the previous backward call
computed, the call for step $t+1$. The two are added before the tanh
derivative is applied:
\begin{lstlisting}
def backward(self, dh_out, dstate_next=None):
    x, h_prev, h_new = self.cache.pop()          # LIFO: this timestep
    H = self.hidden_size
    if dstate_next is None:
        dstate_next = [0.0] * H
    # Two gradient paths into h_t: this step's output, and the future.
    dh_total = [dh_out[i] + dstate_next[i] for i in range(H)]
    # Through tanh: d/dz tanh(z) = 1 - h^2.
    da = [dh_total[i] * (1.0 - h_new[i] * h_new[i]) for i in range(H)]
    # Accumulate weight grads (the += across timesteps), as in Linear.
    for i in range(H):
        for j in range(self.input_size):
            self.dW_xh[i][j] += da[i] * x[j]
        for j in range(H):
            self.dW_hh[i][j] += da[i] * h_prev[j]
        self.db_h[i] += da[i]
    # Backprop to this step's input and to the previous state.
    dx = [sum(self.W_xh[i][j] * da[i] for i in range(H))
          for j in range(self.input_size)]
    dh_prev = [sum(self.W_hh[i][j] * da[i] for i in range(H))
               for j in range(H)]
    return dx, dh_prev
\end{lstlisting}
The single line \texttt{dh\_total = dh\_out + dstate\_next} is the whole of
backpropagation through time: the gradient at step $t$ is what came back from
this step's output plus what came back from the future through $h_{t+1}$.

The weight gradients accumulate across all $T$ timesteps with \texttt{+=}.
This is the same accumulation the library has used since the \texttt{Linear}
layer of \Cref{sec:mlp-backprop}, in a third flavor. Per-example gradients
sum within a mini-batch. Per-position gradients sum within a single
convolutional forward, because the kernel is reused at every spatial
position. Per-timestep gradients sum within one recurrent unroll, because
the recurrent matrix is reused at every step. The same operation runs at a
finer grain each time: across examples, then across positions, then across
time.

Two pieces of bookkeeping make the unroll trainable. The first is gradient
clipping. After all $T$ backward steps, the global L2 norm of the gradient
is capped at $5$, the standard mitigation for the exploding-gradient half of
the problem the next section names. The second is the mean-gradient step.
The gradients summed over the $T$ steps are divided by $T$ before the
update, which turns the sum back into a mean and keeps the effective
learning rate sane regardless of how long the unroll was. This is exactly
the convention \texttt{Network.step(lr, n)} uses for a mini-batch, with $T$
playing the role of the batch size.

A gradient check confirms the derivation. Unroll the cell for $T = 3$ with a
small \texttt{Linear} output projection and \texttt{SoftmaxCrossEntropy},
build the analytic gradient with one forward and backward sweep, and compare
every parameter against a central finite difference. The worst relative
error lands on the order of $10^{-9}$, far under the $10^{-4}$ tolerance the
book has used since \Cref{ch:foundations}. The chapter's paired notebook
regenerates this, and \texttt{scratchnn}'s
\texttt{test\_gradients.py} carries the same check.

\section{Vanishing gradients}\label{sec:vanishing-gradients}

Backpropagation through time is correct, and it is fragile. Each step of the
backward pass multiplies the gradient on the state by a Jacobian:
\begin{equation}\label{eq:rnn-jacobian}
  \frac{\partial h_{t+1}}{\partial h_t}
  = \operatorname{diag}\!\bigl(1 - h_{t+1}^2\bigr)\, W_{hh}.
\end{equation}
The diagonal factor is the tanh derivative, componentwise; the matrix factor
is the recurrent weight matrix. To send a gradient back $k$ steps, the
backward pass multiplies $k$ of these Jacobians together, so the magnitude
of the gradient on an early step is governed by the product. That product is
in turn governed by the spectral properties of $W_{hh}$, scaled by the tanh
derivative, which lies in $(0, 1]$ and is almost always strictly below $1$.

The product of many matrices of roughly constant size grows or shrinks
geometrically. If the spectral radius of the effective Jacobian sits below
$1$, the gradient on early timesteps decays exponentially in $k$, and those
parameters receive essentially no signal from losses many steps later. If it
sits above $1$, the gradient blows up, and one bad sequence can throw the
weights off entirely. This is the vanishing-gradient problem, with exploding
gradients as its twin, and it is the reason a vanilla recurrent network
cannot easily learn dependencies more than a few dozen steps long. The
architecture is the right shape for sequences in principle. The learning
dynamics make long-range information transfer hard in practice.

This is the price of the prior, and the experiment later in this chapter
shows it from both sides. Numerically, the paired notebook pushes a unit
gradient back through the recurrence alone and watches its norm fall by
several orders of magnitude over a few dozen steps, a direct measurement of
\Cref{eq:rnn-jacobian} compounding. Qualitatively, the generated text drifts
over long spans even after the local structure is correct: spelling and
short words come out right while coherence across a sentence wanders. That
drift is not a separate defect. It is the same limit, the state failing to
carry information far enough, showing in the output rather than in the
gradient.

Two mitigations are standard, and the chapter mentions them without building
them. Gated cells, the long short-term memory and the gated recurrent unit,
replace the simple tanh recurrence with an additive update. The long
short-term memory keeps a cell state whose Jacobian is
$\operatorname{diag}(f_t)$ for a learned forget gate $f_t$ with entries in
$[0, 1]$; when the gate sits near $1$, the gradient flows through
unattenuated, a gradient highway across many steps. This library implements
the vanilla \texttt{RNNCell} only, and the extension is mechanical: more
matrix multiplies, one per gate and one for the cell-state update, the same
\texttt{+=} accumulation in the backward pass, and no new conceptual
machinery. Gradient clipping, already in the training loop, handles the
exploding half. The point worth keeping is that this one Jacobian-product
issue is what motivated the entire gated-recurrent literature, and
eventually the architecture the next chapter turns to.
